
\subsection{Delivery conditions from a simulation output}
A person-level Monte Carlo simulation was run under declared assumptions --- two independent risk
draws per couple, classroom-level facts logged once per classroom, a single assessor draw per case,
and pre-arranged versus day-of logistics as one switch --- at 100 classrooms (eight couples each, five
seeds) and again at 1,000 classrooms (three seeds). This is a result of the assumptions coded into
the model, not empirical data about any delivered programme: no participant was observed, and nothing
here is an effectiveness estimate. It is reported because it shows which operator-controllable
conditions the model's own assumptions make consequential, checked by an independent adversarial
re-run across additional seeds.

That re-run confirmed the rankings from the original seeds while showing the first point estimates
were seed-specific and should be reported as ranges. Across seeds 1--8 and both scales, in the
simulation, under its assumptions, pre-arranging an assessor with no stake in the wedding, a private
room, and a professional interpreter \emph{before} the day was the single largest lever: true-risk
cases handled safely rose $\approx\!2.1$--$2.5\times$, cases left stuck with no assessor fell
$2.5$--$4.9\times$, and private slots overheard by family fell $6$--$10\times$. These rankings held
across seeds and at ten-fold scale; the magnitudes did not, hence the ranges.

Two further conditions followed. In the simulation, under its assumptions, once logistics were
pre-arranged, untrained facilitators became the dominant residual cause of unsafe continuation. In the
simulation, under its assumptions, a narrow, trained trigger list reduced false triggers relative to a
broad one. Day length behaved differently: in the simulation, under its assumptions, walk-outs
($\approx$5\% of couples, modelled) were unchanged by any safeguarding lever and moved only when the
dose --- the length of a single day --- was redesigned, for example by moving content to a follow-up
slot. Table~\ref{tab:delivery} summarises the four conditions the output ranked as decisive for
delivery, as distinct from the model elements themselves. None of this is a finding about \tphe{} in
practice; it states what the declared assumptions in the simulation code imply, offered as input to
co-design, not as evidence that any condition improves an outcome.

\begin{table}[H]
\centering\footnotesize\raggedright
\caption{Four delivery conditions the simulation output ranked as decisive (not model elements;
output of assumptions, not data). All figures are simulation outputs of declared assumptions, not
measured effects; the assessor row's ``no stake'' condition is a proposed extension supported only by
analogy evidence (\protect\cite{ref158,ref159,ref162,ref163}; see Results, Table~\protect\ref{tab:rules}).
A case the model counts as ``stuck'' is one where the R4 referral pathway --- to a qualified
multidisciplinary service such as Thailand's One-Stop Crisis Centre (OSCC) or an equivalent
designated service --- could not be completed in the simulation's own assumptions, not one in which
OSCC access itself was tested.}
\label{tab:delivery}
\begin{tabularx}{\columnwidth}{@{}p{2.4cm}Y@{}}
\toprule
\textbf{Condition} & \textbf{What the output attributes to it, in the simulation}\\
\midrule
Assessor, room, and interpreter pre-arranged before the day & Largest lever modelled: safely handled
cases $\times2.1$--$2.5$; stuck cases $\div2.5$--$4.9$; overheard private slots $\div6$--$10$.\\
Facilitator training gate before delivery & Once logistics are pre-arranged, the residual cause of
unsafe continuation in the model.\\
Narrow, trained trigger list & Fewer false triggers than a broad or untrained list, in the model.\\
Follow-up slot rather than a longer single day & The only lever in the model that moved walk-outs,
which safeguarding levers alone did not change.\\
\bottomrule
\end{tabularx}
\end{table}
